\section{Теория поля}
\label{sec:qft-open}

На масштабах ядра и элементарных частиц тот же конус, та же локальность и тот же спин
читаются как лагранжиан с калибровочными симметриями $SU(3)\times SU(2)\times U(1)$,
фермионами поколений и механизмом Хиггса.
Квантовая электродинамика и хромодинамика дают проверенные поправки.

Заряд, который доходит до прибора, кратен элементарному:
\begin{equation}
  Q = n \chargeE, \qquad n\in\spaceZ.
  \label{eq:charge-quantum}
\end{equation}
Свободный кварк не наблюдается, заряд лептона целый.
Закон, у которого устойчивый дефект нёс бы произвольную долю этого кванта,
таблицу частиц не воспроизводит.
Произведение зеркального отражения, зарядового сопряжения и обращения времени опыт сохраняет,
\begin{equation}
  (\mathrm{CPT})\,\mathcal{L}\,(\mathrm{CPT})^{-1} = \mathcal{L};
  \label{eq:cpt}
\end{equation}
по отдельности отражение и зарядовое сопряжение в слабом взаимодействии нарушены.
Общий поворот фазы вакуума $\wavePsi\mapsto e^{i\symAlpha}\wavePsi$ не меняет частот:
это то же сохранение заряда, \eqref{eq:continuity}.
Обе симметрии должны сидеть в самом законе.

Отсюда же видно, где непрерывная запись перестаёт справляться.

Возмущёние в квантовой теории поля содержит расходимости.
Группа ренормализации снимает их с наблюдаемых величин
и оставляет вопрос, что стоит за взаимодействием там,
где оно перестаёт быть точечным.
Это то же ограничение плотности, теперь внутри возмущения.

Стандартная модель содержит более девятнадцати настраиваемых величин:
массы, углы смешивания, константы связи.
Симметрии $SU(3)\times SU(2)\times U(1)$ их не фиксируют.
Как и этажом выше, числа стоят в законе вписанными.

Квантование метрики в той же схеме не замкнуто.
Планковская энергия
\begin{equation}
  E_P = \sqrt{\frac{\planckHbar \speedC^5}{\constGrav}}
  \label{eq:planck-energy}
\end{equation}
--- граница записи тех же требований.
Теория поля --- последний непрерывный этаж.

